\documentclass[11pt,a4paper]{article}
\usepackage[T1]{fontenc}
\usepackage{newtxtext,newtxmath}
\usepackage[left=2.3cm,right=2.3cm,top=2.3cm,bottom=2.3cm]{geometry}
\usepackage{xcolor,tcolorbox,enumitem,titlesec,fancyhdr,url}
\definecolor{brand}{HTML}{1F6F4A}
\definecolor{dark}{HTML}{142B20}
\definecolor{soft}{HTML}{F3F6F1}
\definecolor{tag}{HTML}{E0ECFA}

\titleformat{\section}{\color{brand}\Large\bfseries}{}{0em}{}
\titleformat{\subsection}{\color{dark}\normalsize\bfseries}{}{0em}{}

\pagestyle{fancy}\fancyhf{}
\lhead{\small AISR P04 G11 \textbullet{} Video Scripts}
\rhead{\small Project P04 \textbullet{} Batch B4}
\cfoot{\thepage}

\newtcolorbox{clipbox}[1]{
  colback=soft,colframe=brand,coltitle=white,title={#1},
  boxrule=0.6pt,arc=2pt,fonttitle=\bfseries,
  left=6pt,right=6pt,top=4pt,bottom=4pt,before skip=5pt,after skip=9pt,
}

\setlist[itemize]{leftmargin=1.3em,itemsep=2pt,topsep=2pt}

\begin{document}

\begin{center}
{\color{brand}\Large\bfseries AISR P04 \textbullet{} Group G11 \textbullet{} Video Scripts}\\[4pt]
{\small Agricultural Field Robot for Crop-Row Coverage \textbullet{} CSAI4014}\\[2pt]
{\small Target 20 minutes total \textbullet{} speaking rate $\approx$ 140 wpm}
\end{center}

\vspace{6pt}

\begin{tcolorbox}[colback=tag,colframe=brand,boxrule=0.5pt,arc=2pt]
\textbf{What everyone does}
\begin{itemize}
  \item Each clip = you on camera, your slide full-screen behind, 1080p MP4.
  \item Filename: \texttt{G11\_<slide no>\_<Name>.mp4}, e.g.\ \texttt{G11\_02\_Jahanvi.mp4}.
  \item Say it in your own words. Reading verbatim scores lower (rubric I2).
  \item Send all your clips to Rudra. Rudra stitches them in order and uploads to YouTube as Unlisted.
\end{itemize}
\end{tcolorbox}

\vspace{4pt}

\noindent\textbf{How many clips each person records:}

\begin{center}
\renewcommand{\arraystretch}{1.25}
\begin{tabular}{|l|c|l|}
\hline
\textbf{Member} & \textbf{Clips} & \textbf{Which slides} \\ \hline
Jahanvi Sengar     & 4 & intro, slide 2, slide 3, closing line (slide 16) \\ \hline
Aagman Sharma      & 5 & intro, slides 4, 5, 6, closing line (slide 16) \\ \hline
Rudra Gupta        & 6 & intro, slides 7, 8, 9, slide 12 demo, slide 17 thank-you \\ \hline
Sahil Agarwal      & 4 & intro, slides 10, 11, closing line (slide 16) \\ \hline
Harsh P.\ Sinha    & 5 & intro, slides 13, 14, 15, closing line (slide 16) \\ \hline
\end{tabular}
\end{center}

\vspace{4pt}

% ======================================================================
\section*{\color{brand}Jahanvi Sengar \textbullet{} 4 clips}

\begin{clipbox}{Clip 1 \textbullet{} Intro \textbullet{} 9 s \textbullet{} on slide 1}
``I am Jahanvi Sengar, SAP ID 500120543. I built the field model, the geometry utilities and the Webots world generator.''
\end{clipbox}

\begin{clipbox}{Clip 2 \textbullet{} Slide 2 Problem and Motivation \textbullet{} 1:35}
``A spraying robot has to cover every row of a crop field --- not more, not less. Miss a strip and weeds and disease take over; overlap a strip and you waste chemicals and water. The hard part is that real rows have obstacles in them. Stones from ploughing, irrigation risers and power poles sit right in the lane, and the robot doesn't know where they are in advance. On the picture on the left you can see the field we used: sixteen metres by ten, eleven lettuce rows one metre apart, with six stones and a power pole scattered across the lanes. Existing industrial sprayers --- Thorvald, lane-following tractors --- react to an obstacle by swerving out of the row, which keeps them safe but loses the lane pattern, so part of the crop is simply never treated. That gap is what we are solving: cover every lane, even when something unknown is blocking it.''
\end{clipbox}

\begin{clipbox}{Clip 3 \textbullet{} Slide 3 Objectives \textbullet{} 1:35}
``We set four measurable goals for the project. One: build the field world --- a sixteen-by-ten metre crop field with eleven lettuce rows, stones and poles, driven by a Pioneer 3-AT with lidar, GPS, inertial unit and a one-point-two metre sprayer boom. Two: cover the field row by row using a boustrophedon plan with arc and pivot headland turns, keeping the lane error below five centimetres RMS. Three: avoid the unknown obstacles with zero collisions and at least ninety-eight per cent coverage of the coverable area. And four: compare the method against Bug2 and APF on twenty random fields, and study how turn style, obstacle density, boom width and PID gains change the result. Every claim in this talk will come back to these four goals.''
\end{clipbox}

\begin{clipbox}{Clip 4 \textbullet{} Slide 16 closing line \textbullet{} 8 s}
``Our robot covered 98.85 per cent of the field with zero collisions and kept its lane in all twenty fields.''
\end{clipbox}

% ======================================================================
\section*{\color{brand}Aagman Sharma \textbullet{} 5 clips}

\begin{clipbox}{Clip 1 \textbullet{} Intro \textbullet{} 9 s \textbullet{} on slide 1}
``I am Aagman Sharma, SAP ID 500123462. I built the boustrophedon coverage planner, the figure scripts and the LaTeX report.''
\end{clipbox}

\begin{clipbox}{Clip 2 \textbullet{} Slide 4 Existing Solutions and the Gap \textbullet{} 1:05}
``The table on this slide summarises four classical approaches. Boustrophedon decomposition from Choset is complete and has low overlap, but it assumes obstacles are known in advance --- our field obstacles are not. Artificial potential field is simple and reactive, but it oscillates near narrow gaps and loses the lane. Bug2 is complete for a point robot, but it takes long detour loops. And the industrial approach --- lane following with a certified safety stop under ISO~18497 --- just stops and waits for a human. The gap we are filling is marked at the bottom: a sensor-based method that keeps the lane pattern intact when an unknown obstacle appears, measured against these classical baselines.''
\end{clipbox}

\begin{clipbox}{Clip 3 \textbullet{} Slide 5 Novelty \textbullet{} 1:10}
``There are five things that are new in our work. First, a lane-preserving detour: when a lane segment is blocked we replace exactly that segment with a safe arc that rejoins the same lane, and on twenty random fields we skip zero strips. Second, a robust lidar map for crops: we add ghost-return and tilt filters that removed seventeen phantom obstacles and stopped the robot from getting stuck. Third, one layered agent, one codebase: the same Python agent runs in Webots and in our headless simulator, with coverage within half a per cent. Fourth, a design study: arc turns are seven-point-four per cent faster than pivot turns, and a 1.2 metre boom on 1 metre lanes gives 98.8 per cent coverage with 24 per cent overlap. And fifth, live coverage accounting: a five-centimetre grid, sprayed soil painted blue in Webots, with an on-screen dashboard.''
\end{clipbox}

\begin{clipbox}{Clip 4 \textbullet{} Slide 6 Syllabus Mapping \textbullet{} 1:00}
``This table links the work back to the AISR syllabus. Uninhabited ground robots: the Pioneer 3-AT, its skid-steer, wheel limits and the lidar-GPS-IMU stack. Path planning: our boustrophedon plan, the geometric detour and the Bug2 and APF baselines. Planning to move: the arc and pivot headland turns, pure pursuit and speed scheduling. Application domains: precision agriculture --- spraying, coverage and overlap. Closed-loop control: a PID heading loop tuned from a step response. And perception plus software architecture: lidar clustering and the declarative-procedural-reflexive agent.''
\end{clipbox}

\begin{clipbox}{Clip 5 \textbullet{} Slide 16 closing line \textbullet{} 9 s}
``Compared with Bug2 and APF our method misses four to thirteen times less area and is thirteen to forty-seven per cent faster.''
\end{clipbox}

% ======================================================================
\section*{\color{brand}Rudra Gupta \textbullet{} 6 clips}

\begin{clipbox}{Clip 1 \textbullet{} Intro + welcome \textbullet{} 25 s \textbullet{} on slide 1}
``Good morning. We are Group G11 from Batch B4, and this is our project P04, Agricultural Field Robot for Crop-Row Coverage, for the course CSAI4014. I am Rudra Gupta, SAP ID 500120437. I built the pure-pursuit and PID tracker, the layered agent and the Webots controller, and I merged the repository. Let us introduce the rest of the team and start with the problem.''
\end{clipbox}

\begin{clipbox}{Clip 2 \textbullet{} Slide 7 System Architecture \textbullet{} 1:05}
``The architecture runs every 32 milliseconds --- that is the Webots step. On the left are the three sensors: a 180-ray 2-D lidar, GPS for x and y, and an inertial unit for roll, pitch and yaw. On the right are the three actuators: the four wheel motors in skid steer, the sprayer boom modelled as a Webots Pen, and the dashboard that writes CSV and JSON logs. In the middle sits our layered agent. The declarative layer holds the boustrophedon plan: lanes, turns and spray tags. The procedural layer runs lidar filtering, the obstacle map, the detour splice, and the pure-pursuit plus PID tracker. The reflexive layer is a safety stop: if anything comes within 40 centimetres, the robot stops and backs off. The same file --- \texttt{agri/agent.py} --- runs in Webots and in our headless simulator.''
\end{clipbox}

\begin{clipbox}{Clip 3 \textbullet{} Slide 8 Robot and Environment Model \textbullet{} 1:00}
``This is our robot and world. The Pioneer 3-AT is a four-wheel-drive skid-steer with wheel radius eleven centimetres and a maximum speed of seventy centimetres per second. We treat it as a unicycle, with wheel speeds given by the standard $v$ plus or minus $\omega b / 2$ over $r$ equation. The lidar sweeps 180 degrees with 180 rays at 35 centimetres above the ground. The environment is a 16-by-10 metre crop field with eleven lettuce rows one metre apart, six stones, a power pole, irrigation risers and a post. The world is generated from a seed by \texttt{tools/generate\_world.py}, so every one of our twenty experiments is a different field with the same rules.''
\end{clipbox}

\begin{clipbox}{Clip 4 \textbullet{} Slide 9 Algorithm 1: Coverage Plan + Tracking \textbullet{} 1:10}
``The coverage plan is boustrophedon: lane 1 north, lane 2 south, and so on. Between lanes we choose either an arc turn --- faster, 7.4 per cent --- or a pivot turn, which is slower but needs less space. The figure above shows the two turn shapes side by side. To follow the plan, the robot uses pure pursuit plus PID. Pure pursuit picks a look-ahead point on the planned path; the angle between that point and the current heading is the error $e$. PID gives the angular rate: $\omega = K_p e + K_i \int e + K_d \dot{e}$, with gains $K_p = 2.2$, $K_i = 0.05$, $K_d = 0.12$ --- tuned from a step response. Speed is scheduled: half a metre per second on a lane, 0.3 inside a turn, slowing down below 2 metres clearance; if the heading error is larger than 55 degrees the robot pivots in place.''
\end{clipbox}

\begin{clipbox}{Clip 5 \textbullet{} Slide 12 Webots Demo voice-over \textbullet{} 3:45 \textbullet{} screen capture, no camera}
\textbf{Record a Webots screen capture of \texttt{worlds/crop\_field.wbt} running the proposed method with arc turns. Narrate the following over the recording; pause the simulation where helpful.}\\[4pt]

``This is the Webots world, \texttt{crop\_field.wbt}. The grey ground is soil, the green patches are lettuce rows, and the small grey shapes are the unknown obstacles --- six stones and a power pole. The robot is the Pioneer 3-AT sitting at the bottom-left corner. I'm pressing Run now.

The robot enters lane one going north. The sprayer is on: the ground is being painted blue behind it --- that's our coverage grid. The top-right of the screen is the dashboard: state, coverage per cent, lane number and a minimap. It says `lane-follow'.

The lidar has picked up the first stone. Watch the speed indicator --- it drops from 0.5 to 0.3 because clearance is below two metres. The detour fires: the state changes to `detour', the robot arcs around the stone, and --- here --- it rejoins the same lane. No lane skipped.

End of lane one. The sprayer turns off, the robot takes an arc headland turn, enters lane two, and the sprayer comes back on. That is the boustrophedon pattern.

Here is a pole in lane four. The perception doesn't care that it is thinner than a stone: cluster, fit a circle, confirm over eight scans, detour, rejoin. Coverage on the dashboard has climbed past sixty per cent.

On this lane the robot slows down because two obstacles are close together, and it picks the shorter of the two safe arcs. The reflexive stop layer stayed silent the whole run.

Final coverage: 98.4 per cent in this seed, 414 seconds of wall clock time, zero collisions, RMS lane error around two centimetres. The field is almost entirely blue --- that is the coverage map we log to CSV. That's the method working end to end.''
\end{clipbox}

\begin{clipbox}{Clip 6 \textbullet{} Slide 17 Thank-you + closing line \textbullet{} 25 s}
``As a team we learned a Git workflow, Webots and how to run fair experiments. Our code is on GitHub at \url{github.com/byRudra/AgriCoverageRobot_G11}. Thank you, we are happy to take questions.''
\end{clipbox}

% ======================================================================
\section*{\color{brand}Sahil Agarwal \textbullet{} 4 clips}

\begin{clipbox}{Clip 1 \textbullet{} Intro \textbullet{} 9 s \textbullet{} on slide 1}
``I am Sahil Agarwal, SAP ID 500120211. I built the lidar perception, the lane-preserving detour, the Bug2 and APF baselines, and the slide deck.''
\end{clipbox}

\begin{clipbox}{Clip 2 \textbullet{} Slide 10 Algorithm 2: Lidar Map + Detour \textbullet{} 1:20}
``The lidar hits the lettuce and the soil as well as the stones, so raw returns are noisy. On the left of this slide is the filter pipeline: we drop returns below a height threshold, drop isolated single hits, then cluster the remaining points and fit a circle to each cluster. A cluster is confirmed as an obstacle only after it is seen in eight consecutive scans. That alone removed seventeen phantom obstacles during tuning. On the right is the geometry of the detour. We expand the obstacle by the robot radius plus a 30-centimetre safety margin --- that gives the detour radius $R = r + 0.32 + 0.30$ metres. We cut out the blocked segment of the lane, and we splice in the shorter of two safe arcs so the robot ends up back on the same lane. That is the key point: the plan doesn't change, only one segment does.''
\end{clipbox}

\begin{clipbox}{Clip 3 \textbullet{} Slide 11 Simulation Setup \textbullet{} 1:10}
``The simulator is Webots R2025a with a 32-millisecond step. The controller is Python 3, in a package called \texttt{agri}. Maximum speed is half a metre per second, turn speed is 0.3. The boom is 1.2 metres on 1-metre lane spacing. The detour margin is 30 centimetres with a 5-metre horizon. We run twenty random fields crossed with three methods --- our proposed method, Bug2 and APF. The commands on the right are exactly what you type: open \texttt{worlds/crop\_field.wbt} to see one run; \texttt{generate\_world.py} makes a new seeded field; \texttt{run\_webots.py} drives the batch; and \texttt{run\_experiments.py} plus \texttt{make\_figures.py} produce the statistics and the figures.''
\end{clipbox}

\begin{clipbox}{Clip 4 \textbullet{} Slide 16 closing line \textbullet{} 9 s}
``In industry this suits spraying, weeding and scouting robots, and the coverage map can be used as a treatment record.''
\end{clipbox}

% ======================================================================
\section*{\color{brand}Harsh Prakash Sinha \textbullet{} 5 clips}

\begin{clipbox}{Clip 1 \textbullet{} Intro \textbullet{} 9 s \textbullet{} on slide 1}
``I am Harsh Prakash Sinha, SAP ID 500123860. I built the coverage metrics, the headless simulator and all the experiments and results.''
\end{clipbox}

\begin{clipbox}{Clip 2 \textbullet{} Slide 13 Results 1: Same Field, Three Methods \textbullet{} 0:50}
``On the same seeded field, the three methods look like this. The proposed method covers 98.43 per cent in 414.6 seconds with zero collisions. Bug2 covers only 95.51 per cent in 449 seconds --- it loops around every obstacle, which costs time and leaves gaps. APF covers 97.74 per cent but, in this run, came within 2.1 centimetres of a stone --- it does not reliably clear the obstacle. You can see the three trajectories above: the green one keeps the lane, the orange one takes long Bug2 loops, and the red one wanders.''
\end{clipbox}

\begin{clipbox}{Clip 3 \textbullet{} Slide 14 Results 2: 20 Random Fields \textbullet{} 0:55}
``Across twenty random fields we get four statistics worth remembering. Mean coverage for the proposed method is 98.85 per cent, versus 95.4 for Bug2 and 85.3 for APF. All twenty fields finished with our method --- APF only finished fifteen. Mean time was 412 seconds, versus 472 for Bug2 and 781 for APF. And RMS lane error was 2.1 centimetres, versus 11.3 and 28.3. The box plot above shows the spread: the proposed method is tight and high; the baselines are wide and low. That is what the lane-preserving detour buys you statistically.''
\end{clipbox}

\begin{clipbox}{Clip 4 \textbullet{} Slide 15 Team Contribution \textbullet{} 1:20}
``The table shows who owns what. Jahanvi: the field model and the Webots world generator, and report sections Abstract, 1, 2 and 13. Aagman: the boustrophedon planner with arc and pivot turns, and the figure scripts, and report sections 3 and 4. Rudra: the PID and pure-pursuit tracker, the layered agent and the Webots controller, and sections 5 and 6. Sahil: the lidar filters, the obstacle map, the detour, and the Bug2 and APF baselines, and sections 7 and 8. And me: the coverage metrics, the headless simulator, and the Monte-Carlo and Webots experiments, and sections 9, 12 and the appendix. Every module can be traced to the owner's GitHub commits in the repository.''
\end{clipbox}

\begin{clipbox}{Clip 5 \textbullet{} Slide 16 closing line \textbullet{} 10 s}
``Next steps are a real robot with RTK-GNSS and ROS~2, and obstacles that are not circles, like fences and ditches.''
\end{clipbox}

\end{document}
